\documentclass[10pt,a4paper]{amsart}
\usepackage[margin=19mm]{geometry}
\usepackage{amsmath,amssymb,mathtools}
\usepackage[scr=rsfs,cal=euler]{mathalfa}
\usepackage{xfrac}
\usepackage[unicode,hidelinks,bookmarks=false]{hyperref}
\newcommand{\ie}{i.e.\ }
\newcommand{\cf}{cf.\ }
\newcommand{\resp}{respectively\ }
\newcommand{\Subsection}{\subsection}
\providecommand{\nobreakdash}{\nobreak\hbox{-}}
\setlength{\emergencystretch}{2em}
\multlinegap0pt
\newcommand{\C}{\mathbb{C}}
\newcommand{\Ext}{\operatorname{Ext}}
\newcommand{\Hom}{\operatorname{Hom}}
\newcommand{\I}{\mathcal{I}}
\newcommand{\OO}{\mathcal{O}}
\newcommand{\PP}{\mathbb{P}}
\newcommand{\dual}{\vee}
\newcommand{\vdim}{\operatorname{vdim}}
\theoremstyle{plain}
\newtheorem{theorem}{Theorem}[section]
\newtheorem{proposition}[theorem]{Proposition}
\newtheorem{lemma}[theorem]{Lemma}
\newtheorem{corollary}[theorem]{Corollary}
\theoremstyle{definition}
\newtheorem{definition}[theorem]{Definition}
\theoremstyle{remark}
\newtheorem{remark}[theorem]{Remark}
\title[Planar rank-one sheaves on $\mathbb{P}^3$, obstruction bundles, and divisor-supported Donaldson--Thomas series]{Planar rank-one sheaves on $\mathbb{P}^3$, obstruction bundles, and divisor-supported Donaldson--Thomas series}
\author{Reginald Anderson}
\date{Source: arXiv:2609.28501v1 (CC BY 4.0)}
\begin{document}
\maketitle

\noindent\emph{Source and licence.} The delivered work is the article shown in the title above, version arXiv:2609.28501v1 (CC BY 4.0), distributed under the Creative Commons Attribution 4.0 International licence, as recorded in the arXiv licence metadata for this exact version and shown on the abstract page saved with this item. The full licence notice, which carries the licence URL and the address of that abstract page, is the bundled file \texttt{licenses/arxiv-2609.28501-CC-BY-4.0.txt}.
\par\smallskip

\noindent\textit{Editorial scope.} Counted unit: the article's Proposition (source label prop:tangent-obstruction-fiber) (source0.tex lines 586--606, source label \texttt{prop:tangent-obstruction-fiber}): ``Let be a plane, let , and set . Then there are canonical exact sequences 0 1 P( Z, Z) 1 X(F,F) H 0 ( P(1) ) 0 and 2 X(F,F) 1 P ( Z, Z(1) ). In particular, 1 X(F,F) 2n 3, 2 X(F,F) 2n. Therefore .''. It is delivered together with the article's complete original proof (source0.tex lines 608--647, one proof environment adjacent to the statement, 40 lines), copied line for line from the article's own text. Required context retained uncounted: lem:ext-sequence (lines 455--474); lem:surface-ext-vanishing (lines 540--556). Editorial formatting only: an AMS \texttt{amsart} preamble that restates the article's own macro definitions and its own theorem declarations, and the theorem environments the retained lines use; the article's own class file is neither shipped nor input; the retained environments are renumbered sequentially in order of appearance, whereas the article prints them with its own numbering; the article's equation numbering is not reproduced and every cross-reference inside the retained text resolves to the wrapper's numbering. No citation occurs inside any retained line, so the wrapper carries no bibliography; All retained bodies are exact copies of the bundled \texttt{sources/211592/source0.tex}, so no omission receipt applies. No asset-inclusion command occurs inside any retained span and the package ships no image asset.


\section*{Retained context: lem:ext-sequence (source0.tex lines 455--474)}

\begin{lemma}\label{lem:ext-sequence}
Let $Y$ be a smooth projective threefold, let $i\colon S\hookrightarrow Y$ be a smooth divisor with normal bundle $N_{S/Y}$, and let $E,G$ be coherent sheaves on $S$.
Then there is a functorial exact sequence
\begin{align*}
0 &\to \Ext^1_S(E,G)
\to \Ext^1_Y(i_*E,i_*G)
\to \Hom_S(E,G\otimes N_{S/Y}) \\
&\to \Ext^2_S(E,G)
\to \Ext^2_Y(i_*E,i_*G)
\to \Ext^1_S(E,G\otimes N_{S/Y}) \to 0.
\end{align*}
Moreover, $\Ext^k_S(E,G)=0$ for all $k>2$, and the continuation of the sequence gives a canonical isomorphism
\[
\Ext^3_Y(i_*E,i_*G)\cong \Ext^2_S(E,G\otimes N_{S/Y}).
\]
There is also an Euler characteristic identity
\begin{equation}\label{eq:chi-general}
\chi_Y(i_*E,i_*G)=\chi_S(E,G)-\chi_S(E,G\otimes N_{S/Y}).
\end{equation}
\end{lemma}


\section*{Retained context: lem:surface-ext-vanishing (source0.tex lines 540--556)}

\begin{lemma}\label{lem:surface-ext-vanishing}
Let $P\cong \PP^2$ and let $E=\I_Z$ for $Z\in P^{[n]}$.
Then
\[
\Ext^2_P(E,E)=0,
\qquad
\Ext^2_P(E,E(1))=0,
\qquad
\Hom_P(E,E(1))\cong H^0\bigl(\OO_P(1)\bigr).
\]
Moreover,
\[
\chi_P(E,E)=1-2n,
\qquad
\chi_P(E,E(1))=3-2n.
\]
\end{lemma}


\section{The counted unit: Proposition (source label prop:tangent-obstruction-fiber) and its complete original proof}

\begin{proposition}\label{prop:tangent-obstruction-fiber}
Let $P\subset X$ be a plane, let $Z\in P^{[n]}$, and set $F=\iota_{P*}\I_Z$.
Then there are canonical exact sequences
\begin{equation}\label{eq:tangent-exact-fiber}
0\longrightarrow \Ext^1_P(\I_Z,\I_Z)
\longrightarrow \Ext^1_X(F,F)
\longrightarrow H^0\bigl(\OO_P(1)\bigr)
\longrightarrow 0
\end{equation}
and
\begin{equation}\label{eq:obstruction-fiber}
\Ext^2_X(F,F)\cong \Ext^1_P\bigl(\I_Z,\I_Z(1)\bigr).
\end{equation}
In particular,
\[
\dim \Ext^1_X(F,F)=2n+3,
\qquad
\dim \Ext^2_X(F,F)=2n.
\]
Therefore $\vdim M_n=3$.
\end{proposition}

\begin{proof}
Since $P\subset X=\PP^3$ is a hyperplane divisor, its normal bundle is
\[
N_{P/X}=N_{P/\PP^3}\cong \OO_X(P)|_P\cong\OO_P(1).
\]
Apply Lemma~\ref{lem:ext-sequence} with $Y=X=\PP^3$, $S=P$, and $E=G=\I_Z$.
By Lemma~\ref{lem:surface-ext-vanishing},
\[
\Ext^2_P(\I_Z,\I_Z)=0,
\qquad
\Hom_P(\I_Z,\I_Z(1))\cong H^0\bigl(\OO_P(1)\bigr).
\]
Moreover, $\Ext^3_P(\I_Z,\I_Z)=0$ because $P$ is a smooth projective surface, as recorded in Lemma~\ref{lem:ext-sequence}.
The exact sequence in that lemma therefore gives \eqref{eq:tangent-exact-fiber} and the isomorphism \eqref{eq:obstruction-fiber}.
The same lemma also gives $\chi_P(\I_Z,\I_Z)=1-2n$, while $\Hom_P(\I_Z,\I_Z)\cong\C$.
Together with the vanishing of $\Ext^2_P(\I_Z,\I_Z)$, this implies $\dim\Ext^1_P(\I_Z,\I_Z)=2n$.
This gives $\dim\Ext^1_X(F,F)=2n+3$.

The Euler characteristic identity \eqref{eq:chi-general} and Lemma~\ref{lem:surface-ext-vanishing} imply
\[
\chi_X(F,F)=\chi_P(\I_Z,\I_Z)-\chi_P\bigl(\I_Z,\I_Z(1)\bigr)=-2.
\]
Moreover $F$ is stable, so $\Hom_X(F,F)=\C$.
Since $K_X\cong\OO_X(-4)$, Serre duality on $X=\PP^3$, the projection formula, and the full faithfulness of pushforward along $\iota_P$ give
\[
\begin{aligned}
\Ext^3_X(F,F)
&\cong \Hom_X(F,F\otimes K_X)^\dual \\
&\cong \Hom_X\bigl(\iota_{P*}\I_Z,\iota_{P*}(\I_Z(-4))\bigr)^\dual \\
&\cong \Hom_P\bigl(\I_Z,\I_Z(-4)\bigr)^\dual \\
&\cong H^0\bigl(P,\OO_P(-4)\bigr)^\dual=0.
\end{aligned}
\]
For the last isomorphism, we use $\mathcal Hom_P(\I_Z,\I_Z)\cong\OO_P$.
Hence
\[
\dim \Ext^2_X(F,F)=\dim\Ext^1_X(F,F)-3=2n.
\]
This proves the virtual dimension statement.
\end{proof}

\end{document}
